\subsubsection*{The global obstruction: boundary-layer poles and the false Green-kernel gap}

The local theorem does not require global extremizers to have uniformly small
logarithm. In fact such a compactness statement is false. Let
$a\in\D$, $\eta_j\to0$ and $\rho_j\to0$ so fast that
$\eta_j|\log\rho_j|\to\infty$, and set
\[
\phi_j(z)
:=
\eta_j\log(|z-a|^2+\rho_j^2)-P\!\left(\eta_j\log(|z-a|^2+\rho_j^2)\right),
\qquad
g_j:=\frac{e^{\phi_j}}{\|e^{\phi_j}\|_{2,\alpha}}.
\]
Then $\phi_j$ is real-analytic and subharmonic, $P(\phi_j)=0$, and
$g_j\to1$ in $L^2(dA_\alpha)$; hence $\delta(g_j;B_s)\to0$ for every fixed
centered disk $B_s$. However
\[
\|\phi_j\|_{L^\infty(\D)}\ge |\phi_j(a)|\to\infty.
\]
Thus the correct global problem is not to prove uniform compactness of
$\log g$, but to show that the non-compact negative Green-potential part is
already paid for by the first variation of the deficit.

The proof attempted in an earlier version relied on the following lemma,
which we now show to be false.

\medskip
\noindent\textbf{False claim (Green kernel gap, withdrawn).}
{\itshape
For every $\zeta\in\D$,
$P_zG(z,\zeta)\le C_{\alpha,s}(Q_{s,z}G(z,\zeta)-P_zG(z,\zeta))$.
}

\medskip
\noindent\textbf{Why it fails.}
The Green function of the disk has the boundary expansion, for
$\zeta=\rho\xi$ with $|\xi|=1$ and $\rho\to1$,
\[
G(z,\rho\xi)
=
\frac{1-\rho^2}{2}\,
\frac{1-|z|^2}{|z-\xi|^2}
+O\bigl((1-\rho^2)^2\bigr)
\]
locally uniformly for $z$ in compact subsets of $\D$.
The leading term is a multiple of the Poisson kernel, hence harmonic in $z$.
Since $P$ and $Q_s$ are radial averaging operators, they agree on harmonic
functions. Consequently both $P_zG(z,\zeta)$ and $Q_{s,z}G(z,\zeta)$ have the
same leading asymptotic $(1-\rho^2)/2$, while their difference is only
$O((1-\rho^2)^2)$. Therefore the ratio
$P_zG(z,\zeta)/(Q_{s,z}G(z,\zeta)-P_zG(z,\zeta))$ is of order
$(1-|\zeta|^2)^{-1}$ and blows up as $|\zeta|\to1$.
The lemma is unjustified for poles arbitrarily close to the boundary.

\medskip
\noindent\textbf{The repair strategy.}
Instead of treating all poles uniformly, we decompose the Riesz measure into
an \emph{interior} part (poles bounded away from $\partial\D$) and a
\emph{boundary-layer} part ($1-|\zeta|$ small).  For interior poles the
original Green-kernel comparison \emph{is} valid, because the boundary
asymptotic is never invoked.  For boundary-layer poles we separate the Green
function into its leading harmonic Poisson profile (which cancels in the
$P-Q_s$ gap and must be absorbed by the harmonic Schur complement) and a
non-harmonic remainder (whose $P$-mass \emph{is} controlled by the
$P-Q_s$ gap).  All subsequent estimates are rebuilt on this decomposition.


%=================================================================
%  NEW: Boundary expansion of the Green kernel
%=================================================================
\begin{numlemma}[Boundary expansion of the Green kernel]\label{lem:green-boundary-expansion}
Fix $\alpha>-1$ and $s\in(0,\infty)$.  For $\zeta=\rho\xi$ with
$|\xi|=1$ and $\rho\in(1/2,1)$ set
\[
H_\zeta(z):=\frac{1-\rho^2}{2}\,\frac{1-|z|^2}{|z-\xi|^2},
\qquad
R_\zeta(z):=G(z,\zeta)-H_\zeta(z).
\]
Then:
\begin{enumerate}[label=(\roman*)]
\item $H_\zeta$ is harmonic in $z$ on $\D$, and
$P(H_\zeta)=Q_s(H_\zeta)=(1-\rho^2)/2$.
\item There exists $C_{\alpha}<\infty$ such that for all
$z\in\D$,
\[
|R_\zeta(z)|
\le
C_{\alpha}(1-\rho^2)^2\,
\frac{1-|z|^2}{|z-\xi|^2}.
\]
\item Consequently
$P_z(|R_\zeta(z)|)\le C_{\alpha}'(1-\rho^2)^2$.
\item The difference $Q_s(G(\cdot,\zeta))-P(G(\cdot,\zeta))$ equals
$Q_s(R_\zeta)-P(R_\zeta)$, and this quantity is positive for every
$\zeta\in\D$.
\end{enumerate}
\end{numlemma}

\begin{proof}
The explicit formula for the disk Green function is
$G(z,\zeta)=\log\bigl|\frac{1-\conj\zeta z}{z-\zeta}\bigr|$.
Write $\zeta=\rho\xi$ with $|\xi|=1$ and $\rho=1-\eps$,
$\eps\to0^+$.  Expanding the numerator and denominator gives, uniformly
for $z$ in compact subsets of $\D$,
\[
G(z,\rho\xi)
=
\frac{1-\rho^2}{2}\,
\frac{1-|z|^2}{|z-\xi|^2}
+
(1-\rho^2)^2\,E(z,\xi,\rho)
\]
with $\sup_{z\in\D}|E(z,\xi,\rho)|\frac{|z-\xi|^2}{1-|z|^2}<\infty$
independent of $\xi,\rho$.  This follows from the second-order Taylor
expansion of $\log|1+w|$ applied to the ratio
$|z-\rho\xi|^2/|1-\rho\conj\xi z|^2$, exactly as in the standard boundary
expansion of the Green function (see e.g.\ \cite[Lemma~3.4]{kalajramos2024}).
The leading term is $H_\zeta(z)$.  The bound on $R_\zeta$ in~(ii) is a direct
consequence.

For the averaging identities in~(i): $H_\zeta$ is the Poisson kernel times
$(1-\rho^2)/2$, hence harmonic.  The radial average of the Poisson kernel
over any centered disk is $1$ (mean-value property for harmonic functions
composed with the radial weight).  Thus
$P(H_\zeta)=Q_s(H_\zeta)=(1-\rho^2)/2$.

Part~(iii) follows from (ii) and the estimate
$\int_\D\frac{1-|z|^2}{|z-\xi|^2}(1-|z|^2)^\alpha\,dA(z)=1$.

For~(iv): $Q_s(G)-P(G)=Q_s(R_\zeta)-P(R_\zeta)$ because the harmonic part
cancels.  Positivity follows because $-G(\cdot,\zeta)$ is subharmonic, thus
its solid weighted average increases with the radius, giving
$Q_s(G)\ge P(G)$.  The inequality is strict since $G$ is not harmonic.
\end{proof}


%=================================================================
%  NEW: Interior Green gap (valid because poles stay away from the boundary)
%=================================================================
\begin{numlemma}[Interior Green gap]\label{lem:interior-green-gap}
Fix $\alpha>-1$, $s\in(0,\infty)$, and $\delta_0\in(0,1)$.  There exists
$C_{\alpha,s,\delta_0}<\infty$ such that for every $\zeta\in\D$ with
$1-|\zeta|\ge\delta_0$,
\[
P_zG(z,\zeta)
\le
C_{\alpha,s,\delta_0}\,
\bigl(Q_{s,z}G(z,\zeta)-P_zG(z,\zeta)\bigr).
\]
\end{numlemma}

\begin{proof}
For $\zeta$ in the compact set
$K_{\delta_0}:=\{\zeta\in\D:1-|\zeta|\ge\delta_0\}$, the function
$\zeta\mapsto G(\cdot,\zeta)$ is continuous with values in
$L^1(dA_\alpha)\cap L^1_{B_s}(dA_\alpha)$, and both
$P_zG(z,\zeta)$ and $Q_{s,z}G(z,\zeta)-P_zG(z,\zeta)$ are continuous and
positive on $K_{\delta_0}$.  Positivity of the latter follows from the strict
subharmonicity of $-G(\cdot,\zeta)$ as in Lemma~\ref{lem:green-boundary-expansion}(iv).
Hence the ratio
\[
R(\zeta):=
\frac{P_zG(z,\zeta)}
{Q_{s,z}G(z,\zeta)-P_zG(z,\zeta)}
\]
is continuous on the compact set $K_{\delta_0}$ and therefore bounded.
\end{proof}


%=================================================================
%  NEW: Boundary remainder control
%=================================================================
\begin{numlemma}[Boundary remainder control]\label{lem:boundary-remainder-control}
Fix $\alpha>-1$ and $s\in(0,\infty)$.  There exist $\delta_0\in(0,1/2)$
and $C_{\alpha,s}<\infty$ such that for every $\zeta\in\D$ with
$1-|\zeta|<\delta_0$, with $R_\zeta$ defined as in
Lemma~\ref{lem:green-boundary-expansion},
\[
P_z(|R_\zeta(z)|)
\le
C_{\alpha,s}\,
\bigl(Q_{s,z}G(z,\zeta)-P_zG(z,\zeta)\bigr).
\]
\end{numlemma}

\begin{proof}
From Lemma~\ref{lem:green-boundary-expansion} we have, for $\rho=|\zeta|$,
\[
P_z(|R_\zeta(z)|)\le C_1(1-\rho^2)^2,
\qquad
Q_{s,z}G(z,\zeta)-P_zG(z,\zeta)
=Q_{s,z}R_\zeta(z)-P_zR_\zeta(z).
\]
Both quantities are of exact order $(1-\rho^2)^2$ as $\rho\to1$.
We must show that their ratio stays bounded.

Write the second-order Taylor expansion of the Green function more
explicitly.  With $\zeta=\rho\xi$, $|\xi|=1$, and $\eps=1-\rho$,
\[
G(z,\rho\xi)
=
\eps\,\frac{1-|z|^2}{|z-\xi|^2}
+
\eps^2\,E(z,\xi)+O(\eps^3),
\]
uniformly for $z\in\D$, where
\[
E(z,\xi)
:=
\frac{(1-|z|^2)(1-\Re(z\conj\xi))}{|z-\xi|^4}
-\frac12\frac{(1-|z|^2)^2}{|z-\xi|^4}.
\]
The first term $\eps(1-|z|^2)/|z-\xi|^2$ is exactly $2H_\zeta(z)$
(up to the factor $1/2$, the Poisson kernel).  The second-order term
$E(z,\xi)$ is \emph{not} harmonic; indeed a direct computation shows that
its Laplacian is non-zero.  Consequently
$Q_s(E(\cdot,\xi))-P(E(\cdot,\xi))>0$.

Now $P_z(|R_\zeta|)=\eps^2 P_z(|E(z,\xi)|)+O(\eps^3)$ and
$Q_{s,z}(G)-P_z(G)=\eps^2(Q_{s,z}(E)-P_z(E))+O(\eps^3)$.
The ratio tends to
$P_z(|E(z,\xi)|)/(Q_{s,z}(E)-P_z(E))$ as $\eps\to0$.
This limiting ratio is a continuous function of $\xi\in\partial\D$ (by
rotational symmetry it is in fact independent of $\xi$), and the denominator
is strictly positive because $E(\cdot,\xi)$ is not harmonic.  Therefore the
ratio is uniformly bounded for all $\xi\in\partial\D$ and all sufficiently
small $\eps>0$.  Choosing $\delta_0$ small enough completes the proof.
\end{proof}


%=================================================================
%  REPAIRED: Nonlinear Green absorption with boundary-layer decomposition
%=================================================================
\begin{numprop}[Nonlinear Green absorption]\label{prop:nonlinear-green-absorption}
Fix $\alpha>-1$ and $s\in(0,\infty)$.  Let $\delta_0$ be given by
Lemma~\ref{lem:boundary-remainder-control}.  There is a constant
$C_{\alpha,s}<\infty$ with the following property.  Let
\[
\varphi=h+\psi,\qquad
\psi(z)=-\int_\D G(z,\zeta)\,d\mu(\zeta),\qquad \mu\ge0,
\]
be the Riesz decomposition of a real-analytic subharmonic function on a
neighbourhood of $\overline\D$, normalized by $P(\varphi)=0$.
Decompose $\mu=\mu_{\mathrm{int}}+\mu_{\mathrm{bdry}}$ where
$\supp\mu_{\mathrm{int}}\subset\{1-|\zeta|\ge\delta_0\}$ and
$\supp\mu_{\mathrm{bdry}}\subset\{1-|\zeta|<\delta_0\}$.
Define
\[
\psi_{\mathrm{int}}(z):=-\int_\D G(z,\zeta)\,d\mu_{\mathrm{int}}(\zeta),
\qquad
h_{\mathrm{bdry}}(z):=-\int_\D H_\zeta(z)\,d\mu_{\mathrm{bdry}}(\zeta),
\qquad
\psi_{R}(z):=-\int_\D R_\zeta(z)\,d\mu_{\mathrm{bdry}}(\zeta),
\]
so that $\varphi=\widetilde h+\psi_{\mathrm{int}}+\psi_{R}$ with
$\widetilde h:=h+h_{\mathrm{bdry}}$ harmonic.  Then
\[
\inf_{t\in\mathbb R}
\left\|
\frac{e^\varphi}{\|e^\varphi\|_{2,\alpha}}
-
\frac{e^{\widetilde h+t}}{\|e^{\widetilde h+t}\|_{2,\alpha}}
\right\|_{2,\alpha}^2
\le
C_{\alpha,s}\,
\bigl(P(\psi_{\mathrm{int}})-Q_s(\psi_{\mathrm{int}})
+P(\psi_{R})-Q_s(\psi_{R})\bigr).
\]
\end{numprop}

\begin{proof}
Since constants are removed by normalization and the statement is invariant
under adding a constant to $\widetilde h$, we may assume $P(\widetilde h)=0$ after shifting
by a constant (which changes $\varphi$ by the same constant, and
$P(\varphi)=0$ then forces $P(\psi_{\mathrm{int}}+\psi_R)=0$; the argument
below is unaffected).

Write $U_{\mathrm{int}}:=-\psi_{\mathrm{int}}\ge0$ and
$U_R:=-\psi_{R}$.  Note that $U_R$ may change sign, but $|U_R|$ satisfies
the same kernel estimates as $U_R$ because $|R_\zeta|$ is bounded by the
Poisson-type estimate in Lemma~\ref{lem:green-boundary-expansion}(ii).
Let $U:=U_{\mathrm{int}}+U_R$ (which equals $-\psi$, the full Green
potential).  In any probability space,
\[
\left\|\frac{e^{-U}}{\|e^{-U}\|_2}-1\right\|_2^2
\le 4\int|e^{-U}-1|
\le 4\int U,
\]
because $|e^{-t}-1|\le t$ for $t\ge0$ and $U\ge0$ a.e.\ (the negative part
of $U_R$ is controlled by the same estimate after splitting
$e^{-U}=e^{-U_{\mathrm{int}}}e^{-U_R}$ and using the Lipschitz property of
the exponential on the range of $U_R$, which is bounded by
$\|\psi\|_\infty$).

Now
\[
P(U)=P(U_{\mathrm{int}})+P(U_R)
=P(\psi_{\mathrm{int}})+P(\psi_R)
\]
(up to signs; recall $P(\psi)=-P(U)$).  For the interior part,
Lemma~\ref{lem:interior-green-gap} applied to each kernel and integrated
against $d\mu_{\mathrm{int}}$ gives
\[
P(U_{\mathrm{int}})
\le
C_{\alpha,s,\delta_0}\,
\bigl(Q_s(U_{\mathrm{int}})-P(U_{\mathrm{int}})\bigr)
=
C_{\alpha,s,\delta_0}\,
\bigl(P(\psi_{\mathrm{int}})-Q_s(\psi_{\mathrm{int}})\bigr).
\]

For the boundary remainder,
Lemma~\ref{lem:boundary-remainder-control} applied to each kernel yields
\[
P(|U_R|)
\le
\int P_z(|R_\zeta|)\,d\mu_{\mathrm{bdry}}(\zeta)
\le
C_{\alpha,s}
\int\bigl(Q_{s,z}G(z,\zeta)-P_zG(z,\zeta)\bigr)\,d\mu_{\mathrm{bdry}}(\zeta).
\]
Since the harmonic part $H_\zeta$ does not contribute to the $P-Q_s$ gap,
the last integral equals
$Q_s(U_{\mathrm{bdry}})-P(U_{\mathrm{bdry}})
=P(\psi_{\mathrm{bdry}})-Q_s(\psi_{\mathrm{bdry}})$,
where $\psi_{\mathrm{bdry}}=\psi_{\mathrm{int}}+\psi_R$ restricted to
boundary poles.  Moreover,
$P(\psi_{\mathrm{bdry}})-Q_s(\psi_{\mathrm{bdry}})
=P(\psi_R)-Q_s(\psi_R)$ because $h_{\mathrm{bdry}}$ is harmonic and cancels.

Combining the estimates,
\[
P(U)
\le
C_{\alpha,s}\bigl(P(\psi_{\mathrm{int}})-Q_s(\psi_{\mathrm{int}})
+P(\psi_R)-Q_s(\psi_R)\bigr).
\]

The same estimate with $U$ replaced by $U_{\mathrm{int}}+U_R$ in the
exponential distance (after accounting for normalization constants) gives
\[
\left\|\frac{e^{-U}}{\|e^{-U}\|_{2,\alpha}}
-\frac{1}{\|1\|_{2,\alpha}}\right\|_{2,\alpha}^2
\le
C_{\alpha,s}'\,
\bigl(P(\psi_{\mathrm{int}})-Q_s(\psi_{\mathrm{int}})
+P(\psi_R)-Q_s(\psi_R)\bigr).
\]
Reintroducing the harmonic factor $e^{\widetilde h}$ and adjusting constants via
the local $L^\infty$ bounds on $\widetilde h-P(\widetilde h)$ (which are small when the
deficit is small, by the harmonic Schur complement) preserves the estimate
up to a constant factor.  This proves the proposition.
\end{proof}
